\documentclass[journal]{IEEEtran}
\usepackage[utf8]{inputenc}
\usepackage[T1]{fontenc}
\usepackage[spanish,es-tabla,es-noquoting]{babel}
\usepackage{amsmath,amssymb,amsfonts}
\usepackage{mathtools}
\usepackage{graphicx}
\graphicspath{{../figures/}}
\usepackage{booktabs}
\usepackage{array}
\renewcommand{\arraystretch}{1.15}
\usepackage[font=small,labelfont=bf]{caption}
\usepackage{cite}
\usepackage{hyperref}
\usepackage{microtype}

\hypersetup{%
  hidelinks,
  pdftitle={Ruptura de simetria y especializacion axial en Seeking the Unicorn},
  pdfauthor={Thomas Chisica},
  pdfsubject={Teoria de Grafos 2026-I, Segunda entrega},
  pdfkeywords={producto fuerte, Laplaciano, Fiedler, conductancia, simetria D4, SODCL}
}

\newcommand{\vol}{\mathrm{vol}}
\newcommand{\cut}{\mathrm{cut}}
\newcommand{\Sobs}{S_{\mathrm{obs}}}
\newcommand{\Sing}{S^{\mathrm{ing}}_{\mathrm{F}}}
\newcommand{\Etwo}{\mathcal{E}_2}

\begin{document}

\title{Ruptura de Simetría y Especialización Axial en \emph{Seeking the Unicorn}:\\ una Formulación Espectral del Problema}

\author{Thomas Chisica\\
Universidad del Rosario\\
Teoría de Grafos 2026-I\\
Profesor: Daniel Alfonso Bojacá Torres}

\maketitle

\begin{abstract}
Este trabajo reformula el análisis del paradigma experimental \emph{Seeking the Unicorn} como un problema de teoría espectral de grafos. El tablero, dotado de conectividad de rey, se modela como el producto fuerte $G = P_8 \boxtimes P_8$, y la división espacial del trabajo de cada díada se representa como una partición de $G$ inducida por frecuencias acumuladas de visita. El punto de partida es que la simetría diédrica $D_4$ del tablero cuadrado fuerza la degeneración $\lambda_2 = \lambda_3$ del Laplaciano, de modo que el eigenespacio de Fiedler es bidimensional y una bisección obtenida del signo de un vector numérico arbitrario no constituye una referencia geométrica bien definida. Sobre esta observación se plantea una solución en dos capas: una familia paramétrica de referencias espectrales, sensible a la degeneración y compatible con las direcciones axiales privilegiadas por $D_4$, para comparar particiones humanas tardías; y una señal geométrica temprana, calculada sobre las primeras cinco rondas ausentes comunes de cada díada, para anticipar la estabilización posterior de roles axiales. El manuscrito desarrolla el marco teórico, el modelamiento matemático y el pipeline computacional que soportan ambas capas, y delimita la ruta hacia los resultados de la entrega final.
\end{abstract}

\begin{IEEEkeywords}
teoría de grafos, producto fuerte, Laplaciano, valor de Fiedler, conductancia, ruptura de simetría, coordinación multiagente
\end{IEEEkeywords}

\section{Introducción}

\IEEEPARstart{L}{a} coordinación entre agentes que carecen de comunicación explícita constituye un problema central en sistemas multiagente, ciencia cognitiva y análisis del comportamiento colectivo \cite{mesbahi2010,andrade2021,goldstone2024}. En el paradigma \emph{Seeking the Unicorn}, dos participantes exploran simultáneamente un tablero compartido de $8\times 8$ celdas en busca de un objetivo cuya posición se desconoce. A lo largo de sucesivas rondas, muchas díadas reducen la redundancia de sus visitas y estabilizan convenciones espaciales focales, en particular divisiones de tipo izquierda--derecha o arriba--abajo \cite{andrade2021}. El fenómeno de interés no es entonces la mera existencia de coordinación, sino el hecho de que la geometría del entorno restringe y favorece ciertas formas de organización por encima de otras.

La pregunta que guía este trabajo no es si los humanos ``vencen'' a un algoritmo espectral, sino cómo debe formularse una comparación válida entre las particiones humanas observadas y la estructura espectral del tablero. Una comparación ingenua contra un único vector de Fiedler resulta insuficiente cuando el segundo autovalor no es simple, pues en ese caso el redondeo espectral depende de la base numérica elegida por la descomposición. En el tablero cuadrado del experimento esa dificultad es estructural y no marginal: la simetría del problema induce un eigenespacio de Fiedler degenerado y, con ello, una familia de particiones espectrales geométricamente distintas. Reproducir observaciones ya documentadas por el estudio original \cite{andrade2021} bajo una etiqueta espectral ingenua no constituiría, por tanto, una contribución matemática.

Sobre esta base, el artículo propone una formulación más precisa del problema. En primer lugar, se modela el tablero como el producto fuerte $P_8 \boxtimes P_8$ y se introduce una familia paramétrica de referencias espectrales, sensible a la simetría $D_4$, capaz de distinguir entre un redondeo ingenuo y las direcciones axiales privilegiadas dentro del eigenespacio degenerado. En segundo lugar, se define una señal geométrica temprana, extraída de las primeras rondas sin unicornio comunes a los dos jugadores, para estudiar si la especialización axial posterior puede anticiparse antes de que la división del trabajo se haya consolidado. La combinación de ambas capas permite tratar el fenómeno como un problema de teoría de grafos con una pregunta empírica bien delimitada y con baselines matemáticamente coherentes.

\section{Descripción del problema}

El problema aplicado se origina en el conjunto de datos \emph{Self-Organized Division of Cognitive Labor} (SODCL), asociado al paradigma \emph{Seeking the Unicorn} \cite{andrade2021}. Cada díada juega 60 rondas; en cada una, ambos participantes inspeccionan celdas del tablero sin comunicación explícita y deciden si el unicornio está presente o ausente. La cobertura conjunta, el solapamiento entre jugadores y la estabilidad de las regiones exploradas permiten observar la emergencia de una división del trabajo.

El análisis requiere distinguir con precisión las dos fuentes oficiales de datos, resumidas en la Tabla~\ref{tab:datasets}. El archivo \emph{performances.csv} contiene la trayectoria completa de las 45 díadas a lo largo de las 60 rondas (presentes y ausentes) y se reserva para evaluar la transferencia a medidas de desempeño conductual. El archivo \emph{humans\_only\_absent.csv}, en cambio, conserva únicamente las rondas sin unicornio e incluye variables derivadas ($DLIndex$, \emph{Similarity}, \emph{Consistency}, \emph{Joint}); sobre este subconjunto se formulan todos los cálculos de partición, conductancia y señales geométricas. Esta distinción es metodológicamente importante: las rondas ausentes permiten observar patrones de cobertura sin censura por terminación temprana, mientras que la tabla completa es necesaria para estudiar transferencia hacia desempeño posterior. En consecuencia, el proyecto debe formularse como un reanálisis sobre un subconjunto estructurado del experimento, y no como si cada díada aportara 60 rondas equivalentes a todas las métricas.

\begin{table}[t]
\centering
\caption{Fuentes de datos empleadas en el proyecto.}
\label{tab:datasets}
\begin{tabular}{@{}lrp{0.55\columnwidth}@{}}
\toprule
\textbf{Archivo} & \textbf{Filas} & \textbf{Rol en el análisis} \\
\midrule
\emph{performances.csv} & 5\,400 & 45 díadas, 60 rondas completas; evaluación de transferencia a desempeño. \\
\addlinespace
\emph{humans\_only\_absent.csv} & 1\,244 & Solo rondas ausentes, con métricas derivadas; extracción de señales geométricas y de la partición observada. \\
\bottomrule
\end{tabular}
\end{table}

El problema central puede formularse en dos preguntas complementarias:
\begin{enumerate}
\item \textbf{Pregunta estructural.} ¿Cuál es la referencia espectral correcta para comparar una partición humana cuando el tablero cuadrado produce un eigenespacio de Fiedler degenerado y, por tanto, una bisección de Fiedler no única?
\item \textbf{Pregunta inferencial.} ¿Puede extraerse de las primeras rondas ausentes una señal geométrica capaz de anticipar qué díadas terminarán con especialización axial estable?
\end{enumerate}

Ambas preguntas están conectadas. La primera depende de la geometría del grafo; la segunda, de cómo esa geometría se manifiesta en la dinámica temprana de exploración. El problema no consiste solo en calcular un corte espectral, sino en relacionar una estructura algebraica del tablero con una regularidad conductual observable en los datos.

\section{Objetivo general y objetivos específicos}

\subsection{Objetivo general}

Modelar el experimento \emph{Seeking the Unicorn} mediante teoría de grafos para estudiar cómo la simetría del tablero y la degeneración del eigenespacio de Fiedler se relacionan con la emergencia y la anticipación temprana de especialización espacial axial en díadas humanas.

\subsection{Objetivos específicos}

\begin{enumerate}
\item Modelar el tablero experimental como el producto fuerte $P_8 \boxtimes P_8$ y caracterizar sus propiedades espectrales básicas, en particular la multiplicidad de $\lambda_2$.
\item Distinguir entre una referencia espectral ingenua y una familia paramétrica de referencias sensibles a la simetría diédrica del tablero.
\item Definir una partición observada de cada díada a partir de frecuencias acumuladas de visita en rondas ausentes y evaluarla mediante conductancia.
\item Construir una etiqueta binaria de orientación axial estable a partir de ventanas tardías del experimento.
\item Diseñar una señal geométrica temprana que capture la ruptura inicial de simetría en la exploración del tablero.
\item Comparar esa señal temprana con las métricas oficiales del conjunto de datos mediante un esquema de clasificación reproducible bajo validación cruzada \emph{leave-one-out}.
\item Consolidar una implementación computacional reproducible que sirva de base al desarrollo de resultados y a la entrega final.
\end{enumerate}

\section{Marco teórico}

\subsection{Coordinación y especialización de roles}

El estudio original documentó que la mayoría de díadas humanas logra una división espacial del trabajo y que dicha coordinación se organiza alrededor de unas pocas configuraciones focales \cite{andrade2021}. Trabajos posteriores han situado este paradigma dentro de una literatura más amplia sobre emergencia de roles especializados en grupos pequeños \cite{goldstone2024}, lo que permite interpretar el fenómeno como un caso de \emph{role specialization} y no como una mera descripción empírica de particiones.

\subsection{Producto fuerte y geometría del tablero}

Sea $P_n$ el grafo camino con $n$ vértices. Dados dos grafos $G_1=(V_1,E_1)$ y $G_2=(V_2,E_2)$, su producto fuerte $G_1 \boxtimes G_2$ tiene como conjunto de vértices $V_1\times V_2$; dos vértices $(u_1,u_2)$ y $(v_1,v_2)$ son adyacentes si y solo si en cada coordenada los vértices coinciden o son adyacentes, y al menos una coordenada presenta adyacencia estricta \cite{hammack2011}. La conectividad tipo rey del tablero experimental coincide precisamente con esta construcción, por lo que el espacio de búsqueda se modela de manera natural como
\[
G = P_8 \boxtimes P_8,
\]
con $|V(G)|=64$ y $|E(G)|=210$. Los grados varían entre $3$ en las cuatro esquinas, $5$ en los bordes y $8$ en el interior. Esta representación permite tratar el entorno como un objeto algebraico bien definido, sobre el cual pueden estudiarse cortes, simetrías y propiedades espectrales.

\subsection{Laplaciano, conductancia y referencia de Cheeger}

Sea $G=(V,E)$ un grafo simple no dirigido, con matriz de adyacencia $A$ y matriz diagonal de grados $D$. El Laplaciano combinatorial se define por
\[
L = D - A.
\]
La matriz $L$ es simétrica y semidefinida positiva, y sus autovalores satisfacen
\[
0 = \lambda_1 \le \lambda_2 \le \cdots \le \lambda_n,
\]
con $\lambda_2 > 0$ si y solo si $G$ es conexo \cite{chung1997,vonluxburg2007}. El segundo autovalor $\lambda_2$ se conoce como conectividad algebraica \cite{fiedler1973,mohar1991}; cualquier autovector asociado recibe el nombre de vector de Fiedler e induce la bisección espectral clásica.

Para una partición $(S,\bar S)$, la calidad geométrica se cuantifica mediante la conductancia
\begin{equation}
h(S) = \frac{\cut(S,\bar S)}{\min\bigl(\vol(S),\vol(\bar S)\bigr)},
\label{eq:conductance}
\end{equation}
donde $\cut(S,\bar S)$ denota el número de aristas con un extremo en $S$ y otro en $\bar S$, y $\vol(S)=\sum_{v\in S}\deg(v)$. Como la conductancia se formula en términos de volumen, la referencia isoperimétrica adecuada se expresa mediante el Laplaciano normalizado
\[
\mathcal{L} = D^{-1/2} L D^{-1/2}, \qquad 0 = \mu_1 \le \mu_2 \le \cdots \le \mu_n.
\]
La desigualdad de Cheeger establece la cota
\begin{equation}
\frac{\mu_2}{2} \le h(G) \le \sqrt{2\mu_2},
\label{eq:cheeger}
\end{equation}
que vincula la calidad de los cortes con el espectro del Laplaciano normalizado \cite{chung1997,cheeger1970}. En este trabajo, $L$ se usa para estudiar la geometría del eigenespacio de Fiedler, mientras que $\mathcal{L}$ se reserva para la referencia isoperimétrica. Sobre $G = P_8 \boxtimes P_8$, el corte de Fiedler ingenuo alcanza $h = 28/210 \approx 0{,}1333$, mientras que los cortes axiales LR y TB alcanzan $h = 22/210 \approx 0{,}1048$; en este manuscrito ese valor se usa como baseline \emph{axis-aligned} de referencia dentro de la familia comparada, y no como una afirmación sobre el óptimo global de $h(G)$.

\subsection{Simetría del cuadrado y degeneración del eigenespacio de Fiedler}

El tablero cuadrado posee las simetrías del grupo diédrico $D_4$, generado por la rotación de $90^\circ$ y la reflexión horizontal, con un total de ocho elementos entre rotaciones y reflexiones. Si $P_\sigma$ denota la matriz de permutación asociada a una de estas simetrías sobre $V(G)$, entonces
\[
P_\sigma L = L P_\sigma,
\]
de modo que los eigenespacios de $L$ son invariantes bajo la acción del grupo. En el tablero experimental, la descomposición espectral del Laplaciano produce
\[
\lambda_2 = \lambda_3 \approx 0{,}4164,
\]
y por tanto el eigenespacio de Fiedler
\[
\Etwo = \operatorname{span}\{v_2, v_3\}
\]
tiene dimensión dos. Esta degeneración implica que no existe una única bisección espectral canónica: distintas direcciones dentro de $\Etwo$ pueden inducir cortes geométricamente distintos, y las direcciones axiales canónicas LR y TB son precisamente aquellas compatibles con la acción de $D_4$ sobre $\Etwo$. En consecuencia, una comparación válida entre particiones humanas y cortes espectrales debe incorporar la estructura completa del eigenespacio, y no una elección arbitraria de vector.

\subsection{Métricas informacionales complementarias}

La conductancia resume la calidad geométrica de una partición, pero no cuantifica por sí sola cuán nítida es la asignación funcional de celdas entre los dos jugadores. Por ello, el modelo se complementa con medidas de teoría de la información. La información mutua entre la identidad del jugador y la celda visitada cuantifica cuánto reduce la celda la incertidumbre sobre quién la explora \cite{cover2006}. La divergencia de Jensen--Shannon entre las distribuciones marginales de visita de ambos jugadores mide la separación entre sus territorios explorados \cite{lin1991}. Estas métricas no sustituyen a la conductancia; la complementan, y permiten contrastar la señal espectral con las métricas conductuales $DLIndex$, \emph{Similarity} y \emph{Consistency} ya disponibles en el conjunto de datos oficial.

\section{Modelamiento del problema}

\subsection{Representación del tablero y de las visitas}

A cada celda $(i,j)$ del tablero, con $1 \le i,j \le 8$, se asocia un vértice de $G = P_8 \boxtimes P_8$. Dos celdas son adyacentes si su distancia de Chebyshev es igual a uno, lo que coincide con la conectividad de rey del juego. Para cada díada $d$, jugador $p \in \{1,2\}$, ronda $r$ y celda $v$, se define la variable indicadora $A^{(d)}_{p,r}(v) \in \{0,1\}$, con $A^{(d)}_{p,r}(v) = 1$ si el jugador $p$ visita la celda $v$ en la ronda $r$ y $A^{(d)}_{p,r}(v) = 0$ en caso contrario. Sobre una ventana de rondas $W$, la frecuencia acumulada de visitas viene dada por
\[
f_p^{(d,W)}(v) = \sum_{r \in W} A^{(d)}_{p,r}(v).
\]

\subsection{Partición observada y orientación tardía}

La partición observada de una díada en la ventana $W$ se define por dominancia de frecuencias,
\begin{equation}
\Sobs^{(d,W)} = \bigl\{v \in V : f_1^{(d,W)}(v) > f_2^{(d,W)}(v) \bigr\},
\label{eq:partition}
\end{equation}
y produce una bipartición dura del tablero que puede evaluarse mediante la conductancia \eqref{eq:conductance}. Para estudiar la organización tardía se consideran tres ventanas: $W_1 = [30,60]$, $W_2 = [40,60]$ y $W_3 = [50,60]$, restringidas a rondas ausentes.

Sobre cada ventana se define el campo firmado de dominancia
\[
m_d^{(W)}(v) = f_1^{(d,W)}(v) - f_2^{(d,W)}(v).
\]
Sea $\tilde m_d^{(W)}$ la versión centrada y normalizada de este campo, y sean $a_{\mathrm{LR}}$ y $a_{\mathrm{TB}}$ dos plantillas centradas que representan gradientes izquierda--derecha y arriba--abajo sobre la rejilla. Los puntajes de orientación se definen mediante los productos internos
\[
s_{\mathrm{LR}}^{(d,W)} = \bigl|\langle \tilde m_d^{(W)}, a_{\mathrm{LR}} \rangle\bigr|,
\qquad
s_{\mathrm{TB}}^{(d,W)} = \bigl|\langle \tilde m_d^{(W)}, a_{\mathrm{TB}} \rangle\bigr|.
\]
Una díada se clasifica como \emph{axialmente estable} si las tres ventanas tardías seleccionan la misma orientación no mixta, LR o TB, y como \emph{MIXED} en cualquier otro caso. La variable binaria de interés se define entonces como
\[
y_d = \begin{cases} 1, & \text{si la díada es axialmente estable,}\\ 0, & \text{en caso contrario.}\end{cases}
\]
Esta etiqueta no pretende representar el conjunto completo de estrategias focales documentadas, sino únicamente la especialización axial que la lente espectral puede caracterizar de forma inequívoca.

\subsection{Señal geométrica temprana}

Para cada díada se seleccionan las primeras cinco rondas ausentes comunes a ambos jugadores disponibles en \emph{humans\_only\_absent.csv}, conformando la ventana temprana $W_d^{\mathrm{early}}$. A partir del campo de dominancia normalizado en esa ventana se propone la señal geométrica temprana
\begin{equation}
g_d = \max \bigl\{
\bigl|\langle \tilde m_d^{(W_d^{\mathrm{early}})}, a_{\mathrm{LR}} \rangle\bigr|,\
\bigl|\langle \tilde m_d^{(W_d^{\mathrm{early}})}, a_{\mathrm{TB}} \rangle\bigr|
\bigr\}.
\label{eq:signal}
\end{equation}
La expresión \eqref{eq:signal} cuantifica el grado en que el campo temprano de visitas ya se alinea con alguna de las dos direcciones axiales privilegiadas por la simetría $D_4$ del tablero, antes de que la coordinación entre los jugadores haya cristalizado. Valores altos de $g_d$ indican dominancia espacial compatible con alguno de los ejes privilegiados ya desde la fase inicial del experimento.

\subsection{Conjuntos analíticos y variables de comparación}

El modelamiento distingue dos conjuntos analíticos. El conjunto \emph{all\_dyads\_audit} contiene las 45 díadas y se utiliza para auditoría descriptiva y para la evaluación predictiva completa. El conjunto \emph{primary\_analysis\_set} restringe el análisis a las díadas cuya partición observada $\Sobs$ admite una evaluación válida de conductancia en el sentido de \eqref{eq:conductance}; las 16 díadas excluidas se documentan de manera individual en un reporte de exclusiones, evitando filtrados silenciosos. Esta estratificación hace explícita una limitación del alcance: la representación bipartita del tablero describe bien las estrategias axiales, pero captura con menor resolución categorías como \emph{ALL}, \emph{NOTHING} o \emph{RS}.

Junto con $g_d$, el modelo considera el vector de métricas oficiales tempranas
\[
\mathbf{m}_d = \bigl(DLIndex_d,\ Similarity_d,\ Consistency_d\bigr),
\]
promediado sobre las mismas cinco rondas ausentes comunes. Estas variables permiten comparar la nueva señal geométrica con las métricas oficiales del conjunto de datos.

\section{Solución propuesta}

La solución al problema descrito se organiza en dos capas complementarias que comparten la misma representación del tablero y el mismo conjunto de variables. La primera capa aborda la pregunta estructural mediante una familia paramétrica de baselines espectrales. La segunda capa aborda la pregunta inferencial mediante un esquema de clasificación temprana sobre la etiqueta $y_d$.

\subsection{Componente geométrico-espectral}

La primera parte de la solución consiste en construir una comparación espectral que respete la degeneración del tablero cuadrado. Sea $\Sing$ la bisección de Fiedler ingenua obtenida a partir de un autovector numérico $v_2$,
\begin{equation}
\Sing = \bigl\{v \in V : v_2(v) \ge \operatorname{mediana}(v_2) \bigr\}.
\label{eq:naive}
\end{equation}
Esta referencia se conserva porque representa el redondeo espectral estándar. Sin embargo, para evitar que la comparación dependa de una base arbitraria se introduce la familia paramétrica de cortes
\begin{align}
u_\theta &= \cos(\theta)\,v_2 + \sin(\theta)\,v_3, \label{eq:utheta}\\
S_\theta &= \bigl\{v \in V : u_\theta(v) \ge \operatorname{mediana}(u_\theta) \bigr\},
\end{align}
con $\theta \in [0,\pi)$. Dentro de esta familia aparecen, en particular, las direcciones axiales canónicas LR y TB del eigenespacio degenerado $\Etwo$. La solución propuesta compara entonces cada partición humana con tres referencias: el redondeo ingenuo $\Sing$, la mejor dirección dentro de la familia $\{S_\theta\}_{\theta \in [0,\pi)}$ y los cortes axiales triviales LR/TB. Esta capa no afirma que los participantes humanos optimicen la conductancia; su objetivo es fijar un punto de comparación matemáticamente coherente con la geometría del tablero, de modo que los análisis posteriores no confundan arbitrariedad de la base con diferencia estructural.

\subsection{Componente de especialización temprana}

La segunda parte de la solución aborda la pregunta inferencial. La variable objetivo es la etiqueta binaria $y_d$ de orientación axial estable. Sobre ella se formulan tres modelos de clasificación logística con vectores de características
\begin{align}
X_d^{(1)} &= (g_d),\\
X_d^{(2)} &= (DLIndex_d,\ Similarity_d,\ Consistency_d),\\
X_d^{(3)} &= (g_d,\ DLIndex_d,\ Similarity_d,\ Consistency_d).
\end{align}
El primer modelo evalúa la información puramente geométrica; el segundo toma como referencia las métricas oficiales del conjunto de datos; el tercero examina si la señal propuesta aporta valor complementario a dichas métricas. Los tres modelos se estiman con estandarización, imputación por mediana y validación cruzada \emph{leave-one-out}, reportando el área bajo la curva ROC (AUC) tanto dentro de muestra como bajo validación cruzada. La comparación entre $X_d^{(1)}$, $X_d^{(2)}$ y $X_d^{(3)}$ delimita el valor incremental de la señal geométrica frente a las métricas ya presentes en el conjunto de datos.

\subsection{Pipeline computacional}

La solución completa se organiza en las siguientes etapas:
\begin{enumerate}
\item construir el grafo del tablero y sus operadores espectrales $L$ y $\mathcal{L}$, y verificar numéricamente la degeneración $\lambda_2 = \lambda_3$ y las cotas de Cheeger \eqref{eq:cheeger};
\item calcular el eigenespacio de Fiedler y generar las referencias geométricas $\Sing$, $\{S_\theta\}$, LR y TB;
\item extraer, para cada díada, la partición observada $\Sobs^{(d,W)}$, la orientación axial tardía y la señal temprana $g_d$;
\item medir conductancia, información mutua y divergencia de Jensen--Shannon sobre las particiones observadas;
\item ajustar los tres modelos de clasificación temprana y comparar su desempeño mediante AUC bajo validación cruzada \emph{leave-one-out};
\item contrastar, en una etapa posterior, la orientación axial estable con la tabla completa \emph{performances.csv} para estudiar la transferencia a desempeño conductual.
\end{enumerate}

Esta solución tiene un propósito acotado y preciso. No pretende sustituir las métricas oficiales del experimento por una sola cantidad espectral, ni afirmar que toda coordinación humana relevante puede reducirse a un corte de baja conductancia. Su aporte esperado es otro: ofrecer una comparación geométrica correcta cuando el tablero es espectralmente degenerado y, sobre esa base, construir una señal temprana interpretable de especialización axial, compatible con la literatura previa y abierta a extensiones en trabajos futuros.

\bibliographystyle{IEEEtran}
\begin{thebibliography}{10}

\bibitem{andrade2021}
E. Andrade-Lotero and R. L. Goldstone, ``Self-organized division of cognitive labor,'' \textit{PLOS ONE}, vol. 16, no. 7, p. e0254532, Jul. 2021.

\bibitem{mesbahi2010}
M. Mesbahi and M. Egerstedt, \textit{Graph Theoretic Methods in Multiagent Networks}. Princeton, NJ: Princeton University Press, 2010.

\bibitem{goldstone2024}
R. L. Goldstone, E. Andrade-Lotero, R. D. Hawkins, and M. E. Roberts, ``The emergence of specialized roles within groups,'' \textit{Topics in Cognitive Science}, 2024.

\bibitem{hammack2011}
R. Hammack, W. Imrich, and S. Klav\v{z}ar, \textit{Handbook of Product Graphs}, 2nd ed. Boca Raton, FL: CRC Press, 2011.

\bibitem{fiedler1973}
M. Fiedler, ``Algebraic connectivity of graphs,'' \textit{Czechoslovak Mathematical Journal}, vol. 23, no. 2, pp. 298--305, 1973.

\bibitem{mohar1991}
B. Mohar, ``The Laplacian spectrum of graphs,'' in \textit{Graph Theory, Combinatorics, and Applications}, vol. 2, Y. Alavi et al., Eds. New York: Wiley, 1991, pp. 871--898.

\bibitem{chung1997}
F. R. K. Chung, \textit{Spectral Graph Theory}. Providence, RI: American Mathematical Society, 1997.

\bibitem{vonluxburg2007}
U. von Luxburg, ``A tutorial on spectral clustering,'' \textit{Statistics and Computing}, vol. 17, no. 4, pp. 395--416, Dec. 2007.

\bibitem{cheeger1970}
J. Cheeger, ``A lower bound for the smallest eigenvalue of the Laplacian,'' in \textit{Problems in Analysis}. Princeton, NJ: Princeton University Press, 1970, pp. 195--199.

\bibitem{cover2006}
T. M. Cover and J. A. Thomas, \textit{Elements of Information Theory}, 2nd ed. Hoboken, NJ: Wiley-Interscience, 2006.

\bibitem{lin1991}
J. Lin, ``Divergence measures based on the Shannon entropy,'' \textit{IEEE Transactions on Information Theory}, vol. 37, no. 1, pp. 145--151, Jan. 1991.

\end{thebibliography}

\end{document}
